C'est le chapitre où l'écart entre le cours et la formalisation est le plus large, dans les
deux sens.

Du côté des statistiques, tout est calcul : la moyenne, la variance, l'écart-type sont définis
ici même, et leurs propriétés — linéarité, effet d'une transformation affine, formule de
König--Huygens — se démontrent en développant une somme.

Du côté des probabilités, une loi est une mesure de masse totale $1$, un événement une partie
mesurable, l'espérance une intégrale. Rien de tout cela n'existe au lycée, où les lois
continues sont introduites par leur densité. Le résultat le plus notable du chapitre est
l'espérance de la loi binomiale : admise en première, elle est ici démontrée par une identité
combinatoire — et la bibliothèque ne la contenait pas.

Quatre énoncés ne sont pas traités, et tous pour la même raison de fond. Le théorème de
Moivre--Laplace est un cas particulier du théorème central limite, dont la forme scolaire —
qui parle d'intervalles et de valeurs numériques — demande un travail de traduction à part
entière. Les intervalles $1\sigma$, $2\sigma$, $3\sigma$ réclament un encadrement de la
fonction d'erreur. Et les intervalles de fluctuation et de confiance reposent sur
Moivre--Laplace : ils tombent avec lui.
